\section{Related Work}
\label{sec:related}

\para{The Transformer and Self-Attention} The seminal work by Vaswani et al. \cite{vaswani2017attention} introduced the Transformer architecture, which replaced recurrent and convolutional layers with self-attention as the primary mechanism for sequence modeling. This shift allowed models to capture global dependencies and scale efficiently, leading to the dominance of models like GPT-2 \cite{radford2019language} in natural language processing and Vision Transformers (ViT) \cite{dosovitskiy2021image} in computer vision. Our work builds on this foundation by analyzing the entropy of these attention distributions to measure their informational efficiency.

\para{The Attention Economy} From a socio-cognitive perspective, the "attention economy" describes a system where human attention is treated as a scarce commodity. Pilgrim et al. \cite{pilgrim2021rising} found that media producers generate increasingly higher-entropy content in shorter formats to compete for this limited resource. G{\'o}nzalez de la Torre et al. \cite{gonzalez2024attention} further linked AI-driven engagement to the capture of human attentional patterns through a "technopolitical" lens. We complement these studies by providing a direct mathematical comparison between artificial attention weights and human behavioral trajectories.

\para{Sequential Recommendation} Recommender systems have increasingly adopted self-attention to model sequential user behavior. SASRec \cite{kang2018self} demonstrated that a self-attentive model could outperform RNNs and Markov Chains by dynamically focusing on relevant items in a user's history. Applications in specialized domains, such as fashion \cite{reusable2022fashion}, have further proven the robustness of this architecture. Our research uses SASRec as a representative model for Internet-field attention modeling to validate its structural similarity to generative AI.
